\chapter{Introduction and overview}

%%%%%%%%%%%%%%%%%%%%%%%%%%%%%
\section{Global perspectives}

Quantum computation and quantum information is the study of the information processing using a quantum mechanical system.

\section{Quantum bits}

The state of a quibit is a vector in a two-dimensional complex vector space:
\[
    \ket{\psi} = \alpha \ket{0} + \beta \ket{1}.
\]
The state vector ${\psi}$ is a \emph{linear combination} of the two states $\ket{0}$ and $\ket{1}$,
often called \emph{superposition}.
The states $\ket{0}$ and $\ket{1}$ are called \emph{computational basis states},
and form an orthonormal basis for this vector space.

When we measure a qubit we get either the result \num{0} with probability $|\alpha|^2$ or the result \num{1} with probability $|\beta|^2$.
Obviously, $|\alpha|^2 + |\beta|^2 = 1$ because probabilities must sum up to \num{1}.

The state of a qubit can be interpreted geometrically as a point on the surface of a sphere, called the \emph{Bloch sphere}:
\[
    \ket{\psi} = \cos\left(\frac{\theta}{2}\right) \ket{0} + e^{i\varphi} \sin\left(\frac{\theta}{2}\right) \ket{1}.
\]

Measurement changes the state of a qubit, collapsing it from a superposition to one of the basis states,
the one consistent with the measurement outcome.

\subsection{Multiple qubits}

A two qubit system has four \emph{computational basis states}: $\ket{00}$, $\ket{01}$, $\ket{10}$, and $\ket{11}$.
A pair of qubits can exist in a superposition of these four states:
\[
    \ket{\psi} = \alpha_{00} \ket{00} + \alpha_{01} \ket{01} + \alpha_{10} \ket{10} + \alpha_{11} \ket{11}.
\]
The measurement result $x \in \{0, 1\}^2$ is obtained with probability $|\alpha_{x}|^2$,
and after the measurement the state of the system is $\ket{xy}$.
The \emph{normalization} condition is $\sum_{x \in \{0, 1\}^2} |\alpha_x|^2 = 1$.

It is possible to measure just one of the qubits in a multi-qubit system.
If for a two qubit system a measures gives \num{0} with probability $|\alpha_{00}|^2 + |\alpha_{01}|^2$,
then the post-measurement state of the system is:
\[
    \ket{\psi} = \frac{\alpha_{00} \ket{00} + \alpha_{01} \ket{01}}{\sqrt{|\alpha_{00}|^2 + |\alpha_{01}|^2}}.
\]
The post-measurement state is \emph{re-normalized}, so that it still satisfies the normalization condition.

The \emph{Bell state} or \emph{EPR pair} is a two qubit state that cannot be written as a product of two single qubit states:
\[
    \ket{\Phi^+} = \frac{\ket{00} + \ket{11}}{\sqrt{2}}
\]
The measurement outcomes of the two qubits in a Bell state are \emph{correlated}.
The measurement correlations in the Bell state are stronger than what is possible in a classical system.

%%%%%%%%%%%%%%%%%%%%%%%%%%%%%
\section{Quantum computation}

\subsection{Single qubit gates}

Quantum gates on a single qubit can be described by two-by-two \emph{unitary} matrices.
A matrix is unitary if it satisfies $U^\dagger U = I$.
The unitary constraint is the only constraint on quantum gates,
so that any unitary matrix describes a valid quantum gate.

The NOT gate is specified by the following unitary matrix:
\[
    X = \begin{pmatrix}
        0 & 1 \\
        1 & 0
    \end{pmatrix}.
\]

Two other important single qubit gates are the $Z$ gate
\[
    Z = \begin{pmatrix}
        1 & 0 \\
        0 & -1
    \end{pmatrix},
\]
and the \emph{Hadamard} gate
\[
    H = \frac{1}{\sqrt{2}} \begin{pmatrix}
        1 & 1 \\
        1 & -1
    \end{pmatrix}.
\]

In the Block sphere picture, single qubit gates correspond to rotations and reflections of the sphere.
The Hadamard gate corresponds to a rotation of the sphere about the $\hat{y}$ axis by $\pi/2$ radians,
followed by a rotation about the $\hat{x}$ axis by $\pi$ radians.

The $X$, $Z$, and $H$ gates are represented by the following quantum circuit diagrams:
\[
    \Qcircuit @C=3em @R=0em {
       \lstick{\alpha\ket{0} + \beta\ket{1}} & \gate{X} & \rstick{\beta\ket{0} + \alpha\ket{1}} \qw \\
    }
\]
\[
    \Qcircuit @C=3em @R=.1em {
       \lstick{\alpha\ket{0} + \beta\ket{1}} & \gate{Z} & \rstick{\alpha\ket{0} - \beta\ket{1}} \qw \\
    }
\]
\[
    \Qcircuit @C=3em @R=.1em {
       \lstick{\alpha\ket{0} + \beta\ket{1}} & \gate{H} & \rstick{\frac{\alpha + \beta}{\sqrt{2}}\ket{0} + \frac{\alpha - \beta}{\sqrt{2}}\ket{1}} \qw \\
    }
\]

Since there are infiintely many $2 \times 2$ unitary matrices, there are infinitely many single qubit gates.
An arbitrary $2 \times 2$ unitary matrix can be expressed as:
\[
    U = e^{i\alpha}
    \begin{pmatrix}
        e^{-i\beta/2} & 0 \\
        0 & e^{i\beta/2}
    \end{pmatrix}
    \begin{pmatrix}
        \cos(\gamma/2) & -\sin(\gamma/2) \\
        \sin(\gamma/2) & \cos(\gamma/2)
    \end{pmatrix}
    \begin{pmatrix}
        e^{-i\delta/2} & 0 \\
        0 & e^{i\delta/2}
    \end{pmatrix}
\]

\subsection{Multiple qubit gates}

An arbitrary quantum computation on any number of qubits can be generated by a finite set of gates that is said to be \emph{universal} for quantum computation.

The prototypical two qubit gate is the \emph{controlled-NOT} or CNOT gate. Its matrix representation is
\[
 U_\text{CN} = \begin{pmatrix}
        1 & 0 & 0 & 0 \\
        0 & 1 & 0 & 0 \\
        0 & 0 & 0 & 1 \\
        0 & 0 & 1 & 0
    \end{pmatrix},
\]
while the circuit representation is
\[
    \Qcircuit @C=3em @R=2em {
        \lstick{\ket{A}}& \ctrl{1} & \rstick{\ket{A}} \qw \\
        \lstick{\ket{B}}& \targ & \rstick{\ket{A \oplus B}} \qw
    }
\]
If the control qubit is set to \num{0}, the target qubit is unchanged.
If the control qubit is set to \num{1}, the target qubit is flipped.
The CNOT gate is a generalization of the classical XOR gate: $\ket{A, B} \mapsto \ket{A, A \oplus B}$ 
(where $\oplus$ is addition modulo \num{2}, which is what the XOR gate computes).

Unitary quantum gates are always reversible because $U^\dagger U = I$. 
The inverse of a unitary matrix is also an unitary matrix,
therefore a quantum gate can always be reversed by another quantum gate.

We have the following \emph{universality} result:
\emph{any multi-qubit quantum gate can be expressed as a combination of CNOT gates and single qubit gates.}

\subsection{Measurements in bases other than the computational basis}

The computational basis is just one possible basis for a single qubit.
Another possible choice is the \emph{Hadamard basis}, which consists of the states $\ket{+}$ and $\ket{-}$:
\[
    \ket{+} = \frac{\ket{0} + \ket{1}}{\sqrt{2}}, \quad
    \ket{-} = \frac{\ket{0} - \ket{1}}{\sqrt{2}}.
\]
An arbitrary state $\ket{\psi} = \alpha \ket{0} + \beta \ket{1}$ can be expressed in the Hadamard basis as:
\[
    \ket{\psi} = \frac{\alpha + \beta}{\sqrt{2}} \ket{+} + \frac{\alpha - \beta}{\sqrt{2}} \ket{-}.
\]
It is possible to measure a qubit in the Hadamard basis, 
in which case the measurement result is either ``+'' with probability $|\alpha + \beta|^2/2$ or ``-'' with probability $|\alpha - \beta|^2/2$.

Given any basis state $\ket{a}$ and $\ket{b}$, it is possible to express an arbitrary state $\ket{\psi}$ in the basis $\{\ket{a}, \ket{b}\}$ as $\ket{\psi} = \alpha \ket{a} + \beta \ket{b}$.
If the states are \emph{orthonormal},
it is possible to measure a qubit with respect to the basis $\{\ket{a}, \ket{b}\}$,
giving the measurement result $a$ with probability $|\alpha|^2$ and $b$ with probability $|\beta|^2$.

\subsection{Quantum circuits}

A simple quantum circuit to swap the states of two qubits is given by the following circuit diagram:
\[
    \Qcircuit @C=1em @R=2em {
        \lstick{\ket{A}} & \ctrl{1} & \targ
        & \ctrl{1} & \rstick{\ket{B}} \qw \\
        \lstick{\ket{B}} & \targ & \ctrl{-1}
        & \targ & \rstick{\ket{A}} \qw
    }
\]
Quantum circuits are read left-to-right, 
and each line in the circuit represent a \emph{wire} in the quantum circuit 
(not necessarily a physical wire; represents the passage of time, or a physical particle moving though space).

By convention,
the state input to the circuit is a computational basis state,
usually $\ket{0}$.

In a quantum circuit, loops are not allowed: 
the circuit it is \emph{acyclic}.
Wires can't be joined or split, and each wire can only be acted on by one gate at a time.
Joining wires (FANIN) is not reversible, and splitting wires (FANOUT) is not allowed because of the no-cloning theorem.

The following circuit represents a three qubit controlled-$U$ gate.
\[
    \Qcircuit @C=1em @R1em {
    & \ctrl{1} & \qw \\
    & \multigate{1}{U} & \qw \\
    & \ghost{U} & \qw \\
    }
\]
The gate has one control qubit and two target qubits.

A measruement operation in a quantum circuit converts a single qubit state $\ket{\psi} = \alpha \ket{0} + \beta \ket{1}$ into a classical bit $M$,
which is either \num{0} with probability $|\alpha|^2$ or \num{1} with probability $|\beta|^2$.
\[
    \Qcircuit @C=1em @R=2em {
        \lstick{\ket{\psi}} & \meter & \rstick{M} \cw
    }
\]
A double line wire in a quantum circuit represents a classical bit, while a single line wire represents a qubit.

\subsection{Qubit copying circuit?}

A classical CNOT gate can be used to copy a classical bit $b$ into a second bit that is initialized to \num{0},
since $b \oplus 0 = b$.

It turns out to be impossible to make a copy of an unknown quantum state.
This property is called the \emph{no-cloning theorem}.
It is one of the main differences between classical and quantum information.

\subsection{Bell states}
A Hadamard gate followed by a CNOT gate can be used to create Bell states (or EPR states, or EPR pairs):
\[
    \Qcircuit @C=1em @R=2em {
        & \gate{H} & \ctrl{1} & \qw \\
        & \qw & \targ & \qw
    }
\]
Bell states can be expressed in the computational basis as:
\[
    \ket{\beta_{xy}} = \frac{\ket{0, y} + (-1)^x \ket{1, \bar{y}}}{\sqrt{2}},
\]
where $\hat{y}$ is the bitwise negation of $y$.

Explicitly, the four Bell states are:
\[
    \ket{\beta_{00}} = \frac{\ket{00} + \ket{11}}{\sqrt{2}},\quad
    \ket{\beta_{01}} = \frac{\ket{01} + \ket{10}}{\sqrt{2}}, \quad
\]
\[
    \ket{\beta_{10}} = \frac{\ket{00} - \ket{11}}{\sqrt{2}},
    \quad\ket{\beta_{11}} = \frac{\ket{01} - \ket{10}}{\sqrt{2}}.
\]

\subsection{Quantum teleportation}

\emph{Quantum teleportation} allows to use an entangle EPR pair to send an unknown quantum state $\ket{\psi}$ from Alice to Bob,
using only classical communication and local measurements.

Quantum teleportation can be implemented with the following quantum circuit:
\[
    \Qcircuit @C=1em @R=2em {
        \lstick{\ket{\psi}} & \qw & \ctrl{1} & \gate{H} & \meter & \cw & \rstick{M_1} \cw \\
        \lstick{} & \qw & \targ & \qw & \meter & \rstick{M_2} \cw & \\
        \lstick{} & \qw & \qw & \qw & \qw & \gate{X^{M_2}}  \cwx[-1]  &  \gate{Z^{M_1}}  \cwx[-2]  &  \rstick{\ket{\psi}}  \qw
        \inputgroupv{2}{3}{0.6em}{1.75em}{\ket{\beta_{00}}}  \\ 
    }
\]

It is instructive to analyze the quantum teleportation circuit step by step.
Alice has a qubit in an unknown state $\ket{\psi}=\alpha \ket{0} + \beta \ket{1}$,
and one of the qubits of the EPR pair $\ket{\beta_{00}}$ shared with Bob.
The initial state of the whole system is
\[
    \ket{\psi_0} \otimes \ket{\beta_{00}} = \frac{1}{\sqrt{2}}
    \left[
        \alpha\ket{0} \otimes (\ket{00} + \ket{11}) +
        \beta\ket{1} \otimes (\ket{00} + \ket{11})
    \right].
\]
Alice applies a CNOT gate to her two qubits, so that the state of the system becomes
\[
    \ket{\psi_1} = \frac{1}{\sqrt{2}}
    \left[
        \alpha\ket{0} \otimes (\ket{00} + \ket{11}) + 
        \beta\ket{1} \otimes (\ket{10} + \ket{01})
    \right].
\]
Next, Alice applies a Hadamard gate to her first qubit, so that the state of the system becomes
\[
    \ket{\psi_2} = \frac{1}{2}
    \left[
        \alpha(\ket{0} + \ket{1}) \otimes (\ket{00} + \ket{11}) +
        \beta(\ket{0} - \ket{1}) \otimes (\ket{10} + \ket{01})  
    \right].
\]
By regrouping the terms differently (following their spatial distribution),
we can express the state of the system as
\[
\begin{aligned}
    \ket{\psi_2} = \frac{1}{2} \big[ & \ket{00} \otimes (\alpha\ket{0} + \beta\ket{1}) + 
    \ket{01} \otimes (\alpha\ket{1} + \beta\ket{0}) + \\
    & \ket{10} \otimes (\alpha\ket{0} - \beta\ket{1}) +
    \ket{11} \otimes (\alpha\ket{1} - \beta\ket{0}) 
    \big].
\end{aligned}
\]
From this expression we can see that if Alice measures her two qubits and gets the result $M_1 M_2$,
then Bob's qubit is in the state
\[
    \ket{\psi_B} = \alpha\ket{M_2} + (-1)^{M_1}\beta\ket{\bar{M_2}}.
\]
To know the state of his qubit, Bob needs to know the measurement result $M_1 M_2$ from Alice.
This is what prevents teleportation from being used to transmit information faster than the speed of light; 
the classical communication of the measurement result is limited by the speed of light.
Once Bob knows the measurement result,
he know his state and can therefore apply $Z^{M_1}X^{M_2}$ to his qubit to recover the original state $\ket{\psi}$.
\emph{The time goes from left to right in a quantum diagram, but the matrix products terms on the right happen first.}

Quantum teleportation does not clone the state $\ket{\psi}$;
after Alice's measurement her original qubit is in a basis state $\ket{M_1}$.

%%%%%%%%%%%%%%%%%%%%%%%%%%%%
\section{Quantum algorithms}

\subsection{Classical computation on a quantum computer}
It is possible to simulate a classical logic gate using a quantum circuit.
Any classical circuit can be replaced by an equivalent circuit containing only \emph{reversible} elements,
using the reversible \emph{Toffoli} gate.
\[
    \Qcircuit @C=1em @R=2em {
        \lstick{a} & \ctrl{1} & \rstick{a} \qw \\
        \lstick{b} & \ctrl{1} & \rstick{b} \qw \\
        \lstick{c} & \targ & \rstick{c\oplus ab} \qw
    }
\]
The target bit $c$ is flipped if and only if both control bits $a$ and $b$ are set to \num{1}.
The Toffoli gate is reversible since applying the gate twice returns the original state of the three bits.
The Toffoli gate can be used to simulate NAND gates, and the FANOUT operation.
With these two operations, it is possible to simulate any other element in a classical circuit.

The quantum Toffoli gate is the quantum equivalent of the classical Toffoli gate:
\[
    \Qcircuit @C=1em @R=2em {
        \lstick{\ket{a}} & \ctrl{1} & \rstick{\ket{a}} \qw \\
        \lstick{\ket{b}} & \ctrl{1} & \rstick{\ket{b}} \qw \\
        \lstick{\ket{c}} & \targ & \rstick{\ket{c\oplus ab}} \qw \
    }
\]
The quantum Toffoli gate is represented by an $8 \times 8$ unitary matrix,
and is therefore reversible.
The quantum Toffoli gate can be used to simulate irreversible classical logic gates as the classical Toffoli gate can.
This ensures that a quantum computer can simulate any classical computation.

Non-deterministic classical computers can also be simulated by a quantum computer.
Instead of using the input $\ket{0}$,
a quantum computer can use the input $\ket{+} = H\ket{0}$ to simulate a non-deterministic classical computer.

\subsection{Quantum parallelism}

Suppose we have a classical function $f\colon\{0, 1\}\to \{0, 1\}$.
We denote $U_f$ the map $\ket{x, y} \mapsto \ket{x, y \oplus f(x)}$, which is unitary.
A quantum computer can be used to compute $U_f$:
\[
    \Qcircuit @C=1em @R=2em {
        \lstick{\ket{x}} & \multigate{1}{U_f} & \rstick{\ket{x}} \qw \\
        \lstick{\ket{y}} & \ghost{U_f} & \rstick{\ket{y\oplus f(x)}} \qw
    }
\]
With $\ket{x} = H\ket{0} = (\ket{0} + \ket{1})/\sqrt{2}$ and $\ket{y} = \ket{0}$, we have:
\[
    U_f\left[\frac{\ket{0} + \ket{1}}{\sqrt{2}} \otimes \ket{0} \right] = \frac{1}{\sqrt{2}} \left( \ket{0, f(0)} + \ket{1, f(1)} \right).
\]
We have evaluated the function $f$ on both inputs \num{0} and \num{1} simultaneously,
using a single evaluation of $U_f$.
This is referred to as \emph{quantum parallelism}.

Quantum parallelism is not immediately useful,
since measuring the output will produce either $f(0)$ or $f(1)$.
To be useful,
quantum computation also requires the ability to extract information about more than one value of $f$ from a superposition.

Using the \emph{Hadamard transform} (or \emph{Walsh-Hadamard transform}) $H^{\otimes n}$ on $n$ qubits,
 it is possible to create a superposition of all $2^n$ possible inputs to the function $f$:
 \[
   H^{\otimes n}\ket{0}^{\otimes n} = \frac{1}{\sqrt{2^n}} \sum_{x \in \{0, 1\}^n} \ket{x}.
 \]
Therefore, the Hadamard transform produces an equal superposision of all computational basis states of $n$ qubits.

Quantum parallelism can be generalized to functions $f\colon\{0, 1\}^n \to \{0, 1\}$ using the Hadamard transform.
Prepare the system in the state $\ket{0}^{\otimes n} \otimes \ket{0}$, and apply $H^{\otimes n}$ to the first $n$ qubits.
Then, apply the unitary map $U_f$ to the system, which produces the state
\[
    \frac{1}{\sqrt{2^n}} \sum_{x \in \{0, 1\}^n} \ket{x} \otimes \ket{f(x)}.
\]

\subsection{Deutsch's algorithm}
Deutsch's algorithm shows how quantum circuits can outperform classical circuits.
Deutsch's algorithm combines quantum parallelism with \emph{interference}.

We consider the following quantum circuit:
\[
    \Qcircuit @C=1em @R=2em {
        \lstick{\ket{0}} & \gate{H} & \multigate{1}{U_f} & \gate{H} & \qw \\
        \lstick{\ket{1}} & \gate{H} & \ghost{U_f} &  \qw & \qw
    }
\]
The input of the quantum circuit is $\ket{\psi_0} = \ket{0}\otimes\ket{1}$.
 After the two Hadamard gates the system is in state
 \[
    \ket{\psi_1} = \frac{\ket{0} + \ket{1}}{\sqrt{2}} \otimes \frac{\ket{0} - \ket{1}}{\sqrt{2}}.
 \]
Applying the unitary map $U_f$ to the system produces the state
% distinguish the cases $f(0) = f(1)$ and $f(0) \neq f(1)$
% use braces for the two cases
\[
    \ket{\psi_2} =
    \begin{cases}
        \pm\frac{\ket{0} + \ket{1}}{\sqrt{2}} \otimes \frac{\ket{0} - \ket{1}}{\sqrt{2}},& \text{if } f(0) = f(1) \\
        \pm\frac{\ket{0} - \ket{1}}{\sqrt{2}} \otimes \frac{\ket{0} - \ket{1}}{\sqrt{2}},& \text{if } f(0) \neq f(1)
    \end{cases}
\]
The final Hadamard gate only acts on the firs qubit, and produces the state
\[
    \ket{\psi_3} =
    \begin{cases}
        \pm\ket{0} \otimes \frac{\ket{0} - \ket{1}}{\sqrt{2}},& \text{if } f(0) = f(1) \\
        \pm\ket{1} \otimes \frac{\ket{0} - \ket{1}}{\sqrt{2}},& \text{if } f(0) \neq f(1)
    \end{cases}
\]
Since we have $f(0)\oplus f(1) = 0$ if and only if $f(0) = f(1)$, we can write
\[
    \ket{\psi_3} = \pm\ket{f(0)\oplus f(1)} \otimes \frac{\ket{0} - \ket{1}}{\sqrt{2}}.
\]
By measuring the first qubit, we can determine $f(0)\oplus f(1)$,
which is a global property of the function $f$,
with a single evaluation $U_f$.
A classical circuit would require two evaluations of $f$ to determine $f(0)\oplus f(1)$.

\subsection{The Deutsch-Jozsa algorithm}
The Deutsch-Jozsa algorithm is a generalization of Deutsch's algorithm to functions $f\colon\{0, 1\}^n \to \{0, 1\}$.

Alice selects a number $x \in \{0, 1\}^n$ and sends it to Bob.
Bob computes $f(x)$ and sends the result back to Alice.
Bob uses either a constant function, or a balanced function (0 for half the inputs, 1 for the other half).
Alice has to determine whether Bob's function is constant or balanced,
given the result of the measurement.

The best classical algorithm requires $2^n/2 + 1$ evaluations of $f$ in the worst case.
Each time has to send Bob $n$ bits of information.

The Deutsch-Jozsa algorithm can determine the nature of the function with a single evaluation of $f$.
The quantum circuit for the Deutsch-Jozsa algorithm is the following:
\[
    \Qcircuit @C=1em @R=2em {
        \lstick{\ket{0}^{\otimes n}} & \gate{H^{\otimes n}} & \multigate{1}{U_f} & \gate{H^{\otimes n}} & \qw \\
        \lstick{\ket{1}} & \gate{H} & \ghost{U_f} &  \qw & \qw
    }
\]
The input state of the circuit is $\ket{\psi_0} = \ket{0}^{\otimes n} \otimes \ket{1}$.
By applying the Hadamard transform to the first $n$ qubits and the Hadamard gate to the last qubit, we get
\[
    \ket{\psi_1} = \frac{1}{\sqrt{2^n}} \sum_{x \in \{0, 1\}^n} \ket{x} \otimes \frac{\ket{0} - \ket{1}}{\sqrt{2}}.
\]
The query register is in an equal superposition of all possible inputs to the function $f$,
while the answer register is in an evenly weighted superposition of the two possible outputs of $f$.
Bob evaluates the function $f$ using the unitary map $U_f: \ket{x}\otimes\ket{y} \mapsto \ket{x}\otimes\ket{y \oplus f(x)}$,
which produces the state
\[
    \ket{\psi_2} = \frac{1}{\sqrt{2^n}} \sum_{x \in \{0, 1\}^n} (-1)^{f(x)} \ket{x} \otimes \frac{\ket{0} - \ket{1}}{\sqrt{2}}.
\]
The function evaluation is now stored in the amplitude of the superposition state of the query register.

The effect of the Hadamard transform on the state $\ket{x}$ is given by
\[
    H^{\otimes n}\ket{x_1}\otimes\cdots\otimes\ket{x_n} = \frac{1}{\sqrt{2^n}} \sum_{z \in \{0, 1\}^n} (-1)^{x\cdot z} \ket{z},
\]
where $x\cdot z$ is the bitwise inner product of $x$ and $z$.

Applying the Hadamard transform to the query register produces the state
\[
    \ket{\psi_3} = \frac{1}{2^n} \sum_{z \in \{0, 1\}^n} \sum_{x \in \{0, 1\}^n} (-1)^{x\cdot z + f(x)}\ket{z} \otimes \frac{\ket{0} - \ket{1}}{\sqrt{2}}.
\]
The amplitude of the state $\ket{0}^{\otimes n}$ is $\sum_x (-1)^{f(x)}/2^n$.

If $f$ is a constant function, then the amplitude of $\ket{0}^{\otimes n}$ is $\pm 1$.
Since $\ket{\psi_3}$ is normalized, 
it follows that the amplitude of all other amplitudes is \num{0},
and a measurement of the query register will always produce the result $\ket{0}^{\otimes n}$.

If $f$ is a balanced function,
then the positive and negative amplitudes for $\ket{0}^{\otimes n}$ cancel out,
and the amplitude of $\ket{0}^{\otimes n}$ is \num{0}.
Therefore, the measurement of the query register will never produce the result $\ket{0}^{\otimes n}$.

If alice measures all \num{0}s then the function is constant; otherwise, the function is balanced.

\subsection{Quantum algorithms summarized}
There are three main classes of quantum algorithms providing an advantage over known classical algorithms:
\begin{itemize}
    \item algorithms based on the \emph{quantum Fourier transform},
    \item \emph{quantum search} alforithms, and
    \item \emph{quantum simulation} algorithms.
\end{itemize}
